\subsection{Properties relative to finite extensions of the ring of scalars}

In this no., A and B denote two integrally closed Noetherian domains such that \(A \subset B\) and B is a finitely generated A-module and K and L the fields of fractions of A and B respectively. We shall write \(div_{A}, \chi_{A}, c_{A}, \gamma_{A}, r_{A}\) instead of div, \(\chi, c, \gamma, r\) respectively where A-modules are concerned and use analogous notation for B-modules.

We know (§ 1, no. 10) that for a prime ideal P of B to be of height 1, it is necessary and sufficient that p = P ∩ A be of height 1; moreover (loc. cit., Proposition 14) for p ∈ P(A), there is only a finite number of prime ideals P ∈ P(B) lying above p.~To abbreviate, we shall denote by P\textbar p the relation ``P lies above p'' (that is p = P ∩ A); we shall then denote by e\_\{P/p\} or e(P/p) the ramification index e(v\_\{P\}/v\_\{p\}) of the valuation v\_\{P\} over the valuation v\_\{p\} (Chapter VI, § 8, no. 1) and by f\_\{P/p\} or f(P/p) the residue class degree f(v\_\{P\}/v\_\{p\}) (loc. cit.); recall that the discrete valuations v\_\{p\} and v\_\{P\} are normed and that f\_\{P/p\} is the degree of the field of fractions of B/P over the field of fractions of A/p.~We set n = r\_A(B), where B is considered as an A-module; hence by definition n = {[}L: K{]} and, for all p ∈ P(A), n is also the rank of the free A\_p-module B for all P\textbar p.~Then it follows from Chapter VI, § 8, no. 5, Theorem 2 that for all p ∈ P(A):

\[
\sum_ {\mathfrak {P} / \mathfrak {p}} e _ {\mathfrak {P} / \mathfrak {p}} f _ {\mathfrak {P} / \mathfrak {p}} = n.\tag{9}
\]

This being so, as D(A) and D(B) are free Z-modules, we define an increasing homomorphism of ordered groups N: D(B) → D(A) (also denoted by \(N_{B/A}\) ), by the condition:

\[
\mathrm{N} (\mathfrak {P}) = f _ {\mathfrak {P} / \mathfrak {p}}. \mathfrak {p} \quad \text {   for   } \quad \mathfrak {P} \in \mathrm{P} (\mathrm{B}), \quad \text {   where   } \quad \mathfrak {p} = \mathfrak {P} \cap \mathrm{A}.\tag{10}
\]

On the other hand (§ 1, no. 10, Proposition 14) we have defined an increasing homomorphism of ordered groups \(i\colon\mathbf{D}(\mathbf{A})\to\mathbf{D}(\mathbf{B})\) (also denoted by \(i_{\mathbf{B}/\mathbf{A}}\) ), by the condition:

\begin{enumerate}
\def\labelenumi{(\arabic{enumi})}
\setcounter{enumi}{10}
\tightlist
\item
\end{enumerate}

\[
i (\mathfrak {p}) = \sum_ {\mathfrak {P} / \mathfrak {p}} e _ {\mathfrak {P} / \mathfrak {p}}. \mathfrak {P} \quad \text {   for   } \mathfrak {p} \in \mathrm{P} (\mathrm{A}).
\]

Clearly for every family \((d_{\iota})\) (resp. \((d_{\iota}^{\prime}))\) of divisors of A (resp. B):

\begin{enumerate}
\def\labelenumi{(\arabic{enumi})}
\setcounter{enumi}{11}
\tightlist
\item
\end{enumerate}

\[
\mathrm{N} (\sup (d _ {\iota} ^ {\prime})) = \sup (\mathrm{N} (d _ {\iota} ^ {\prime})), \quad \mathrm{N} (\inf (d _ {\iota} ^ {\prime})) = \inf (\mathrm{N} (d _ {\iota} ^ {\prime}))\tag{13}
\]

\[
i \left(\sup (d _ {\iota})\right) = \sup (i (d _ {\iota})), \quad i (\inf (d _ {\iota})) = \inf (i (d _ {\iota})).
\]

Formula (9) shows that:

\[
\mathbf {N} \circ i = n. 1 _ {\mathbf {D} (\mathbf {A})}.\tag{14}
\]

For all \(a \in \mathbf{A}\) ( \(\S 1\) , no. 10, Proposition 14):

\[
i \left(\operatorname{div} _ {\mathbf {A}} (a)\right) = \operatorname{div} _ {\mathbf {B}} (a).\tag{15}
\]

We deduce (by means of (13)) that, for every fractional ideal a of A, also:

\[
i \left(\operatorname{div} _ {\mathbf {A}} (\mathfrak {a})\right) = \operatorname{div} _ {\mathbf {B}} (\mathfrak {a B}).\tag{16}
\]

For every element \(b \in \mathbf{B}\) , we know (Chapter V, §1, no. 3, Corollary to Proposition 11) that \(\mathrm{N}_{\mathrm{L / K}}(b) \in \mathbf{A}\) ; moreover (Chapter VI, §8, no. 5, formula (9)):

\[
v _ {\mathfrak {p}} \left(\mathrm{N} _ {\mathrm{L} / \mathbf {K}} (b)\right) = \sum_ {\mathfrak {P} / \mathfrak {p}} f _ {\mathfrak {P} / \mathfrak {p}} v _ {\mathfrak {P}} (b)\tag{17}
\]

whence:

\[
\mathbf {N} (\operatorname{div} _ {\mathbf {B}} (b)) = \operatorname{div} _ {\mathbf {A}} (\mathbf {N} _ {\mathbf {L} / \mathbf {K}} (b)).\tag{18}
\]

Formulae (15) and (18) show that, by taking quotients, the homomorphisms N and i define homomorphisms which will also be denoted, by an abuse of language, by:

\[
\mathrm{N}: \mathrm{C} (\mathrm{B}) \rightarrow \mathrm{C} (\mathrm{A}), \quad i: \mathrm{C} (\mathrm{A}) \rightarrow \mathrm{C} (\mathrm{B}).
\]

Note that the homomorphism \(i\colon\mathbf{C}(\mathbf{A})\to\mathbf{C}(\mathbf{B})\) is not in general injective (§ 3, Exercise 7).

Recall that for every B-module R, \(R_{[A]}\) denotes the A-module obtained from R by restricting the scalars to A (Algebra, Chapter II, §1, no. 13).

PROPOSITION 18. (i) For R to be a pseudo-zero B-module, it is necessary and sufficient that the A-module \(\mathbf{R}_{[\mathrm{A}]}\) be pseudo-zero.

\begin{enumerate}
\def\labelenumi{(\roman{enumi})}
\setcounter{enumi}{1}
\tightlist
\item
  For R to be finitely generated torsion B-module, it is necessary and sufficient that \(R_{[A]}\) be a finitely generated torsion A-module and then:
\end{enumerate}

\[
\chi_ {\mathbf {A}} (\mathbf {R} _ {[ \mathbf {A} ]}) = \mathbf {N} (\chi_ {\mathbf {B}} (\mathbf {R})).\tag{19}
\]

\begin{enumerate}
\def\labelenumi{(\roman{enumi})}
\setcounter{enumi}{2}
\tightlist
\item
  For R to be finitely generated B-module, it is necessary and sufficient that \(R_{[A]}\) be a finitely generated A-module and then:
\end{enumerate}

\begin{enumerate}
\def\labelenumi{(\arabic{enumi})}
\setcounter{enumi}{19}
\tightlist
\item
\end{enumerate}

\[
c _ {\mathbf {A}} (\mathrm{R} _ {[ \mathbf {A} ]}) = \mathrm{N} (c _ {\mathbf {B}} (\mathrm{R})) + r _ {\mathbf {B}} (\mathrm{R}) c _ {\mathbf {A}} (\mathrm{B})\tag{21}
\]

\[
r _ {\mathrm{A}} (\mathrm{R} _ {[ \mathrm{A} ]}) = n. r _ {\mathrm{B}} (\mathrm{R}) \quad (\text { recall   that } n = r _ {\mathrm{A}} (\mathrm{B})).
\]

As B is a finitely generated A-module, for R to be a finitely generated B-module, it is necessary and sufficient that \(R_{[A]}\) be a finitely generated A-module. Moreover, if b is the annihilator of R, \(b \cap A = a\) is the annihilator of \(R_{[A]}\) ; as B is integral over A, there is no ideal other than 0 lying above the ideal 0 of A (Chapter V, §2, no. 1, Corollary 1 to Proposition 1) and hence it amounts to the same to say that \(\mathfrak{a} \neq 0\) or that \(\mathfrak{b} \neq 0\) .

\begin{enumerate}
\def\labelenumi{(\roman{enumi})}
\item
  By virtue of this last remark, we may confine our attention to the case where R is a torsion B-module. If b is contained in a prime ideal \(\mathfrak{P} \in \mathrm{P}(B)\) , a is contained in \(\mathfrak{P} \cap A = p\) , which is of height 1. Conversely, if a is contained in a prime ideal \(p \in \mathrm{P}(A)\) , there exists a prime ideal P of B which contains b and lies above p (Chapter V, §2, no. 1, Corollary 2 to Theorem 1). Assertion (i) follows from these remarks and no. 4, Definition 2.
\item
  For every finitely generated torsion B-module R, we write
\end{enumerate}

\[
\phi (\mathbf {R}) = \chi_ {\mathbf {A}} (\mathbf {R} _ {[ \mathbf {A} ]});
\]

clearly (for finitely generated torsion B-modules) \(\phi\) satisfies conditions (1) and (2) of Proposition 11 of no. 5 (taking account of (i)). There therefore exists a homomorphism \(\theta: \mathrm{D}(\mathrm{B}) \to \mathrm{D}(\mathrm{A})\) such that \(\phi(\mathrm{R}) = \theta(\chi_{\mathbf{B}}(\mathrm{R}))\) for every finitely generated torsion B-module R. The homomorphism \(\theta\) is determined by its value for every B-module of the form \(\mathrm{B}/\mathfrak{P}\) where \(\mathfrak{P} \in \mathrm{P}(\mathrm{B})\) , since \(\chi_{\mathbf{B}}(\mathrm{B}/\mathfrak{P}) = \mathfrak{P}\) . Now, for every prime ideal \(\mathfrak{q} \neq \mathfrak{p} = \mathfrak{P} \cap \mathrm{A}\) in \(\mathrm{P}(\mathrm{A})\) , \(\mathfrak{p} \notin \mathfrak{q}\) and hence \((\mathrm{B}/\mathfrak{P})_{\mathfrak{q}} = 0\) . On the other hand, if we set \(\mathrm{S} = \mathrm{A} - \mathfrak{p}\) , \(\mathfrak{P} \cdot \mathrm{S}^{-1}\mathrm{B}\) is a maximal ideal of \(\mathrm{S}^{-1}\mathrm{B}\) and \((\mathrm{B}/\mathfrak{P})_{\mathfrak{p}} = \mathrm{S}^{-1}\mathrm{B}/\mathfrak{P}\) . \(\mathrm{S}^{-1}\mathrm{B}\) is isomorphic to the field of fractions of \(\mathrm{B}/\mathfrak{P}\) (Chapter II, § 2, no. 5, Proposition 11), that is to the residue field of \(v_{\mathfrak{P}}\) ; its length as an \(\mathrm{A}_{\mathfrak{p}}\) -module is therefore \(f_{\mathfrak{P}/\mathfrak{p}}\) ; which proves that \(\theta = \mathrm{N}\) (no. 5, Definition 4).

\begin{enumerate}
\def\labelenumi{(\roman{enumi})}
\setcounter{enumi}{2}
\tightlist
\item
  If T is the torsion submodule of R, \(T_{[A]}\) is the torsion submodule of \(R_{[A]}\) and \((\mathrm{R}/\mathrm{T})_{[\mathrm{A}]} = \mathrm{R}_{[\mathrm{A}]}/\mathrm{T}_{[\mathrm{A}]}\) ; to prove (21) we may therefore confine our attention to the case where R is torsion-free. Then R is identified with a sub-B-module of \(\mathrm{R}_{(\mathrm{L})}\) and contains a basis \((e_{i})_{1 \leqslant i \leqslant m}\) of \(\mathrm{R}_{(\mathrm{L})}\) over L. If \((b_{j})_{1 \leqslant j \leqslant n}\) is a basis of L over K consisting of elements of B, the \(b_{j}e_{i}\) form a basis of \(\mathrm{R}_{(\mathrm{L})}\) over K consisting of elements of R, whence (21). On the other hand let M be a free sub-B-module of R such that R/M is a torsion B-module; as \(M_{[A]}\) is a direct sum of \(r_{\mathrm{B}}(\mathrm{R})\) A-modules isomorphic to B, by Proposition 16 (i)
\end{enumerate}

\[
c _ {\mathbf {A}} (\mathbf {M} _ {[ \mathbf {A} ]}) = r _ {\mathbf {B}} (\mathbf {R}). c _ {\mathbf {A}} (\mathbf {B}).
\]

Moreover, \(c_{\mathbf{A}}((\mathbf{R} / \mathbf{M})_{[\mathbf{A}]})=-c_{\mathbf{A}}(\mathbf{N}(\chi_{\mathbf{B}}(\mathbf{R} / \mathbf{M})))\) by virtue of (19); but by definition of the homomorphism \(\mathrm{N}\colon\mathrm{C}(\mathrm{B})\to\mathrm{C}(\mathrm{A})\) , \(c_{\mathbf{A}}(\mathrm{N}(d))=\mathrm{N}(c_{\mathbf{B}}(d))\) for all \(d\in\mathrm{D}(\mathrm{B})\) and, as \(c_{\mathbf{B}}(\chi(\mathrm{R}/\mathrm{M}))=-c_{\mathbf{B}}(\mathrm{R})\) by definition, finally

\[
c _ {\mathbf {A}} ((\mathbf {R} / \mathbf {M}) _ {[ \mathbf {A} ]}) = \mathbf {N} (c _ {\mathbf {B}} (\mathbf {R}));
\]

then it suffices to apply Proposition 16 (i) to obtain (20).

PROPOSITION 19. Let R be a finitely generated B-module. For R to be reflexive, it is necessary and sufficient that \(R_{[A]}\) be a reflexive A-module.

We have remarked in the proof of Proposition 18 that for R to be a torsion-free B-module, it is necessary and sufficient that \(R_{[A]}\) be a torsion-free A-module. We may therefore assume that R is a lattice of \(W = R \otimes_{B} L\) with respect to B. We shall use the following lemma:

LEMMA 4. Let W be a vector space of finite rank over L and let R be a lattice of W with respect to B. Then, for all \(\mathfrak{p} \in \mathrm{P}(\mathbf{A})\) , \((\mathbf{R}_{[\mathbf{A}]})_{\mathfrak{p}} = \bigcap_{\mathfrak{P}/\mathfrak{p}} \mathbf{R}_{\mathfrak{P}}\) .

If \(S = A - p\) , the prime ideals of the ring \(S^{-1}B\) are generated by the prime ideals of \(B\) not meeting \(S\) , in other words the ideals \(\mathfrak{P}_i\) ( \(1 \leqslant i \leqslant m\) ) lying above \(p\) and the ideal (0); this shows that \(S^{-1}B\) is a semi-local ring whose maximal ideals are the \(m_i = \mathfrak{P}_i(S^{-1}B)\) for \(1 \leqslant i \leqslant m\) ; moreover the local ring \((S^{-1}B)_{\mathfrak{P}_i}\) is isomorphic to \(B_{m_i}\) (Chapter II, § 2, no. 5, Proposition 11) and hence is a discrete valuation ring. The ring \(S^{-1}B\) is therefore a Dedekind domain (§ 2, no. 2, Theorem 1 (f)) and, as it is semi-local, it is a principal ideal domain (§ 2, no. 2, Proposition 1). This being so, \((R_{[A]})_p\) is equal to \(S^{-1}R\) considered as an \(A_p\) -module; by the above, \(S^{-1}R\) is a free lattice of \(W\) with respect to \(S^{-1}B\) and Theorem 2 of no. 2 may therefore be applied to it, giving \(S^{-1}R = \bigcap_i (S^{-1}R)_{m_i}\) ; but \((S^{-1}R)_{m_i} = R_{\mathfrak{P}_i}\) , which proves the lemma.

Returning to the proof of Proposition 19, by Lemma 4,

\[
\bigcap_ {\mathfrak {P} \in P (B)} R _ {\mathfrak {P}} = \bigcap_ {\mathfrak {p} \in P (A)} \left(R _ {[ A ]}\right) _ {\mathfrak {p}}
\]

and the conclusion follows from no. 2, Theorem 2.

COROLLARY. The ring B is a reflexive A-module.

PROPOSITION 20. (i) For a finitely generated A-module M to be pseudo-zero, it is necessary and sufficient that \(\mathbf{M} \otimes_{\mathbf{A}} \mathbf{B}\) be a pseudo-zero B-module.

\begin{enumerate}
\def\labelenumi{(\roman{enumi})}
\setcounter{enumi}{1}
\tightlist
\item
  If \(\mathbf{M}\) is a finitely generated torsion A-module, \(\mathbf{M} \otimes_{\mathbf{A}} \mathbf{B}\) is a finitely generated B-module and:
\end{enumerate}

\[
\chi_ {\mathbf {B}} (\mathbf {M} \otimes_ {\mathbf {A}} \mathbf {B}) = i (\chi_ {\mathbf {A}} (\mathbf {M})).\tag{22}
\]

\begin{enumerate}
\def\labelenumi{(\roman{enumi})}
\setcounter{enumi}{2}
\tightlist
\item
  If \(\mathbf{M}\) is a finitely generated A-module, \(\mathbf{M} \otimes_{\mathbf{A}} \mathbf{B}\) is a finitely generated B-module and:
\end{enumerate}

\begin{enumerate}
\def\labelenumi{(\arabic{enumi})}
\setcounter{enumi}{22}
\tightlist
\item
\end{enumerate}

\[
c _ {\mathbf {B}} (\mathbf {M} \otimes_ {\mathbf {A}} \mathbf {B}) = i (c _ {\mathbf {A}} (\mathbf {M}))\tag{24}
\]

\[
r _ {\mathbf {B}} (\mathbf {M} \otimes_ {\mathbf {A}} \mathbf {B}) = r _ {\mathbf {A}} (\mathbf {M}).
\]

\begin{enumerate}
\def\labelenumi{(\roman{enumi})}
\tightlist
\item
  Let \(\mathfrak{P}\) be a prime ideal of \(\mathbf{B},\mathfrak{p} = \mathfrak{P}\cap \mathbf{A}\) ; then \((\mathbf{M}\otimes_{\mathbf{A}}\mathbf{B})_{\mathfrak{P}} = \mathbf{M}\otimes_{\mathbf{A}}\mathbf{B}_{\mathfrak{P}}\) (Chapter II, § 2, no. 7, Proposition 18) and on the other hand
\end{enumerate}

\[
\mathbf {M} \otimes_ {\mathbf {A}} \mathbf {B} _ {\mathfrak {p}} = (\mathbf {M} \otimes_ {\mathbf {A}} \mathbf {A} _ {p}) \otimes_ {\mathbf {A} _ {p}} \mathbf {B} _ {\mathfrak {p}} = \mathbf {M} _ {p} \otimes_ {\mathbf {A} _ {p}} \mathbf {B} _ {\mathfrak {p}};
\]

the relation \(\mathbf{M}_{\mathfrak{p}} = 0\) is therefore equivalent to \((\mathbf{M} \otimes_{\mathbf{A}} \mathbf{B})_{\mathfrak{P}} = 0\) (Chapter II, §4, no. 4, Lemma 4). It suffices to apply this remark to the ideal \(\mathfrak{P} = (0)\) and the ideals \(\mathfrak{P} \in \mathbf{P}(\mathbf{B})\) to prove (i), taking account of no. 4, Definition 2.

To prove (ii), we shall use the following lemma:

LEMMA 5. Let \(\mathbf{M}_1, \mathbf{M}_2\) be two finitely generated A-modules, and \(f: \mathbf{M}_1 \to \mathbf{M}_2\) an injective homomorphism. Then the kernel of \(f \otimes 1_{\mathbf{B}}: \mathbf{M}_1 \otimes_{\mathbf{A}} \mathbf{B} \to \mathbf{M}_2 \otimes_{\mathbf{A}} \mathbf{B}\) is pseudo-zero.

Let \(\mathfrak{p}\) be a prime ideal of A of height \(\leqslant 1\) . Then \((\mathbf{M}_i\otimes_{\mathbf{A}}\mathbf{B})_{\mathfrak{p}} = (\mathbf{M}_i)_{\mathfrak{p}}\otimes_{\mathbf{A}_{\mathfrak{p}}}\mathbf{B}_{\mathfrak{p}}\) \((i = 1,2)\) (Chapter II, §2, no. 7, Proposition 18) and \((f\otimes 1_{\mathbf{B}})_{\mathfrak{p}} = f_{\mathfrak{p}}\otimes 1_{\mathbf{B}_{\mathfrak{p}}}\) ; the hypothesis that \(f\) is injective implies that so is \(f_{\mathfrak{p}}\) (Chapter II, §2, no. 4, Theorem 1); on the other hand, in view of the choice of \(\mathfrak{p},\mathbf{A}_{\mathfrak{p}}\) is a principal ideal domain and \(\mathbf{B}_{\mathfrak{p}}\) a finitely generated torsion-free \(\mathbf{A}_{\mathfrak{p}}\) -module and hence free; we conclude that \(f_{\mathfrak{p}}\otimes 1_{\mathbf{B}_{\mathfrak{p}}}\) is itself injective. If \(\mathbf{I} = \operatorname {Ker}(f\otimes 1)\) , then \(\mathbf{I}_{\mathfrak{p}} = \operatorname {Ker}((f\otimes 1)_{\mathfrak{p}})\) (Chapter II, §2, no. 4, Theorem 1); therefore \(\mathrm{I_p} = 0\) , whence a fortiori \(\mathrm{I}_{\mathfrak{P}} = (\mathrm{I}_{\mathfrak{p}})_{\mathfrak{P}} = 0\) for \(\mathfrak{P}| \mathfrak{p}\) , which proves the lemma (no. 4, Definition 2).

We return now to the proof of (ii). For every finitely generated torsion A-module M, write \(\phi(M) = \chi_{B}(M \otimes_{A} B)\) ; it follows from (i) that, if M is pseudo-zero, \(\phi(M) = 0\) . On the other hand, consider an exact sequence of finitely generated torsion A-modules:

\[
0 \rightarrow \mathrm{M} _ {1} \rightarrow \mathrm{M} _ {2} \rightarrow \mathrm{M} _ {3} \rightarrow 0.
\]

It follows from Lemma 4 that there is an exact sequence of B-modules:

\[
0 \rightarrow \mathrm{I} \rightarrow \mathrm{M} _ {1} \otimes_ {\mathrm{A}} \mathrm{B} \rightarrow \mathrm{M} _ {2} \otimes_ {\mathrm{A}} \mathrm{B} \rightarrow \mathrm{M} _ {3} \otimes_ {\mathrm{A}} \mathrm{B} \rightarrow 0
\]

where I is pseudo-zero. Using no. 5, Corollary to Proposition 10, we therefore obtain \(\phi(\mathbf{M}_2) = \phi(\mathbf{M}_1) + \phi(\mathbf{M}_3)\) . We therefore conclude from Proposition 11 of no. 5 that there exists a homomorphism \(\theta: \mathbf{D}(\mathbf{A}) \to \mathbf{D}(\mathbf{B})\) such that \(\phi(\mathbf{M}) = \theta(\chi_{\mathbf{A}}(\mathbf{M}))\) for every finitely generated torsion A-module M. To prove that \(\theta = i\) , it suffices to show that \(\phi(\mathbf{A}/\mathfrak{p}) = i(\mathfrak{p})\) for all \(\mathfrak{p} \in \mathbf{P}(\mathbf{A})\) ; now \((\mathbf{A}/\mathfrak{p}) \otimes_{\mathbf{A}} \mathbf{B} = \mathbf{B}/\mathfrak{p}\mathbf{B}\) and, for all \(\mathfrak{P} \in \mathbf{P}(\mathbf{B})\) , \((\mathbf{B}/\mathfrak{p}\mathbf{B})_{\mathfrak{P}} = \mathbf{B}_{\mathfrak{P}}/\mathfrak{p}\mathbf{B}_{\mathfrak{P}}\) ; the last module is 0 if \(\mathfrak{P}\) does not lie above \(\mathfrak{p}\) ; if on the contrary \(\mathfrak{P}| \mathfrak{p}, \mathbf{B}_{\mathfrak{P}}/\mathfrak{p}\mathbf{B}_{\mathfrak{P}}\) is a \(\mathbf{B}_{\mathfrak{P}}\) -module of length \(e(\mathfrak{P}/\mathfrak{p})\) by definition of the ramification index (Chapter VI, § 8, no. 1);

therefore \(\chi_{\mathbf{B}}(\mathrm{B} / p\mathrm{B}) = \sum_{\mathfrak{P} / p}e_{\mathfrak{P} / p}\cdot \mathfrak{P} = i(p)\) , which completes the proof of (ii).

Formula (24) is immediate, for

\[
(\mathbf {M} \otimes_ {\mathbf {A}} \mathbf {B}) \otimes_ {\mathbf {B}} \mathbf {L} = \mathbf {M} \otimes_ {\mathbf {A}} \mathbf {L} = (\mathbf {M} \otimes_ {\mathbf {A}} \mathbf {K}) \otimes_ {\mathbf {K}} \mathbf {L}
\]

and the rank of \((\mathbf{M} \otimes_{\mathbf{A}} \mathbf{K}) \otimes_{\mathbf{K}} \mathbf{L}\) over L is equal to the rank of \(M \otimes_{A} K\) over K. To show (23), consider a free submodule H of M such that Q = M/H is a torsion A-module. Applying Lemma 4 as above, we obtain an exact sequence of B-modules:

\[
0 \rightarrow \mathrm{I} \rightarrow \mathrm{H} \otimes_ {\mathrm{A}} \mathrm{B} \rightarrow \mathrm{M} \otimes_ {\mathrm{A}} \mathrm{B} \rightarrow \mathrm{Q} \otimes_ {\mathrm{A}} \mathrm{B} \rightarrow 0
\]

where I is pseudo-zero. It therefore follows from no. 7, Proposition 16 (ii) and (v) and Corollary 1 to Proposition 16 that

\[
c _ {\mathrm{B}} (\mathrm{M} \otimes_ {\mathrm{A}} \mathrm{B}) = c _ {\mathrm{B}} (\mathrm{Q} \otimes_ {\mathrm{A}} \mathrm{B}) = - c _ {\mathrm{B}} \left(\chi_ {\mathrm{B}} (\mathrm{Q} \otimes_ {\mathrm{A}} \mathrm{B})\right) = - c _ {\mathrm{B}} (i \left(\chi_ {\mathrm{A}} (\mathrm{Q})\right))
\]

by virtue of (ii); but by definition of the homomorphism \(i: \mathbf{C}(\mathbf{A}) \to \mathbf{C}(\mathbf{B})\) , \(c_{\mathbf{B}}(i(\chi_{\mathbf{A}}(\mathbf{Q}))) = i(c_{\mathbf{A}}(\chi_{\mathbf{A}}(\mathbf{Q}))) = -i(c_{\mathbf{A}}(\mathbf{M}))\) , which completes the proof of (23).

\textbf{Remarks}\par

\begin{enumerate}
\def\labelenumi{(\arabic{enumi})}
\item
  If M is a reflexive A-module, \(M \otimes_{A} B\) is not necessarily reflexive (exercise 6). However it is so when B is a flat A-module (no. 2, Proposition 8).
\item
  Let C be a third integrally closed Noetherian domain such that \(B \subset C\) and C is a finitely generated B-module (and hence also a finitely generated A-module). Then we have the transitivity equations:
\item
\end{enumerate}

\[
\mathbf {N} _ {\mathbf {C} / \mathbf {A}} = \mathbf {N} _ {\mathbf {B} / \mathbf {A}} \circ \mathbf {N} _ {\mathbf {C} / \mathbf {B}},\tag{26}
\]

\[
i _ {\mathbf {C} / \mathbf {A}} = i _ {\mathbf {C} / \mathbf {B}} \circ i _ {\mathbf {B} / \mathbf {A}}
\]

which follow immediately from the transitivity equations for the ramification indices and the residue class degrees (Chapter VI, § 8, no. 1, Lemma 1).

